% !TeX root = ../main.tex
% ============================================================
% LÁMINA 19 — TRABAJOS FUTUROS (fondo claro)
% Datos: latex/capitulos/07_Conclusiones.tex (§2, Trabajos futuros)
% ============================================================
\begin{frame}[plain]
  \begin{tikzpicture}[remember picture, overlay]

    \fill[paperBg] (current page.south west) rectangle (current page.north east);

    \node[anchor=north west, align=left]
      at ([xshift=1.4cm, yshift=-0.55cm]current page.north west)
      {\kicker[accentBlue]{18 \,$\cdot$\, Trabajos futuros}};

    \node[anchor=north west, align=left, text width=13.2cm]
      at ([xshift=1.4cm, yshift=-1.05cm]current page.north west)
      {\Large\bfseries\sffamily\color{inkText}Líneas de continuación tras las limitaciones detectadas};

    % --- 4 tarjetas claras en fila ---
    \def\fW{3.075cm}
    \def\fGap{0.3cm}
    \def\fH{5.20cm}
    \def\fTextW{2.28cm}
    \def\fY{-2.05cm}

    % ---------- Tarjeta 1 — Verificación entre etapas ----------
    \node[anchor=north west, fill=cardWhite, draw=lineGrey, line width=0.6pt,
          rounded corners=6pt, minimum width=\fW, minimum height=\fH, inner sep=0pt]
      (f1) at ([xshift=1.4cm, yshift=\fY]current page.north west) {};
    \fill[accentTeal, rounded corners=2pt]
      (f1.north west) rectangle ([yshift=-0.12cm]f1.north east);
    \node[anchor=north west, circle, draw=accentTeal, line width=1pt, minimum size=16pt, inner sep=0pt]
      at ([xshift=0.4cm, yshift=-0.42cm]f1.north west)
      {\fontsize{7.5}{7.5}\selectfont\bfseries\sffamily\color{accentTeal}1};
    \node[anchor=north west, text width=\fTextW, align=left] (f1title)
      at ([xshift=0.4cm, yshift=-0.95cm]f1.north west)
      {\small\bfseries\sffamily\color{inkText}Verificación entre etapas};
    \node[anchor=north west]
      at ([yshift=-0.18cm]f1title.south west)
      {\parbox[t]{\fTextW}{\raggedright\fontsize{7.5}{9}\selectfont\sffamily\color{mutedText}%
       Validar entidades \textit{NER} contra el texto fuente y confirmar negaciones antes de
       normalizar.}};

    % ---------- Tarjeta 2 — Mejora del NER ----------
    \node[anchor=north west, fill=cardWhite, draw=lineGrey, line width=0.6pt,
          rounded corners=6pt, minimum width=\fW, minimum height=\fH, inner sep=0pt]
      (f2) at ([xshift=1.4cm+\fW+\fGap, yshift=\fY]current page.north west) {};
    \fill[accentTeal, rounded corners=2pt]
      (f2.north west) rectangle ([yshift=-0.12cm]f2.north east);
    \node[anchor=north west, circle, draw=accentTeal, line width=1pt, minimum size=16pt, inner sep=0pt]
      at ([xshift=0.4cm, yshift=-0.42cm]f2.north west)
      {\fontsize{7.5}{7.5}\selectfont\bfseries\sffamily\color{accentTeal}2};
    \node[anchor=north west, text width=\fTextW, align=left] (f2title)
      at ([xshift=0.4cm, yshift=-0.95cm]f2.north west)
      {\small\bfseries\sffamily\color{inkText}Mejora del \textit{NER}};
    \node[anchor=north west]
      at ([yshift=-0.18cm]f2title.south west)
      {\parbox[t]{\fTextW}{\raggedright\fontsize{7.5}{9}\selectfont\sffamily\color{mutedText}%
       Diccionario de siglas y abreviaturas clínicas; mejor delimitación de \textit{spans}
       compuestos para reducir la sobre-segmentación.}};

    % ---------- Tarjeta 3 — Normalización con contexto ----------
    \node[anchor=north west, fill=cardWhite, draw=lineGrey, line width=0.6pt,
          rounded corners=6pt, minimum width=\fW, minimum height=\fH, inner sep=0pt]
      (f3) at ([xshift=1.4cm+2*\fW+2*\fGap, yshift=\fY]current page.north west) {};
    \fill[accentTeal, rounded corners=2pt]
      (f3.north west) rectangle ([yshift=-0.12cm]f3.north east);
    \node[anchor=north west, circle, draw=accentTeal, line width=1pt, minimum size=16pt, inner sep=0pt]
      at ([xshift=0.4cm, yshift=-0.42cm]f3.north west)
      {\fontsize{7.5}{7.5}\selectfont\bfseries\sffamily\color{accentTeal}3};
    \node[anchor=north west, text width=\fTextW, align=left] (f3title)
      at ([xshift=0.4cm, yshift=-0.95cm]f3.north west)
      {\small\bfseries\sffamily\color{inkText}Normalización con contexto};
    \node[anchor=north west]
      at ([yshift=-0.18cm]f3title.south west)
      {\parbox[t]{\fTextW}{\raggedright\fontsize{7.5}{9}\selectfont\sffamily\color{mutedText}%
       \textit{Entity linking} sobre el dictado completo, no por coincidencia aislada; mejoraría
       la baja cobertura de \textit{CIE-10-ES}.}};

    % ---------- Tarjeta 4 — Menos alucinaciones del LLM ----------
    \node[anchor=north west, fill=cardWhite, draw=lineGrey, line width=0.6pt,
          rounded corners=6pt, minimum width=\fW, minimum height=\fH, inner sep=0pt]
      (f4) at ([xshift=1.4cm+3*\fW+3*\fGap, yshift=\fY]current page.north west) {};
    \fill[accentTeal, rounded corners=2pt]
      (f4.north west) rectangle ([yshift=-0.12cm]f4.north east);
    \node[anchor=north west, circle, draw=accentTeal, line width=1pt, minimum size=16pt, inner sep=0pt]
      at ([xshift=0.4cm, yshift=-0.42cm]f4.north west)
      {\fontsize{7.5}{7.5}\selectfont\bfseries\sffamily\color{accentTeal}4};
    \node[anchor=north west, text width=\fTextW, align=left] (f4title)
      at ([xshift=0.4cm, yshift=-0.95cm]f4.north west)
      {\small\bfseries\sffamily\color{inkText}Menos alucinaciones del \textit{LLM}};
    \node[anchor=north west]
      at ([yshift=-0.18cm]f4title.south west)
      {\parbox[t]{\fTextW}{\raggedright\fontsize{7.5}{9}\selectfont\sffamily\color{mutedText}%
       \textit{RAG} o verificación \textit{post-generación} que contraste cada afirmación contra
       las entidades y el texto fuente.}};

    % --- Nota inferior ---
    \draw[lineGrey, opacity=0.6, line width=0.5pt]
      ([yshift=-0.3cm]f1.south west) -- ([yshift=-0.3cm]f4.south east);

    \node[anchor=north west]
      at ([xshift=1.4cm, yshift=-7.25cm]current page.north west)
      {\parbox[t]{13.2cm}{\raggedright\footnotesize\sffamily\color{mutedText}%
       Además, evaluar el \textit{pipeline} sobre \textit{hardware} con mayor \textit{VRAM}
       permitiría usar \textit{LLMs} de mayor escala que \texttt{llama3.1:8b} y aislar si las
       limitaciones observadas dependen del tamaño del modelo o de la arquitectura secuencial.}};

  \end{tikzpicture}
  \end{frame}
